\documentclass[UTF8,fontset=fandol,11pt]{ctexart}
\usepackage[a4paper,margin=24mm]{geometry}
\usepackage{amsmath,amssymb,amsthm,graphicx,tabularx,booktabs,hyperref}
\newtheorem{definition}{定义}[section]
\newtheorem{theorem}[definition]{定理}
\newtheorem{example}[definition]{例题}
\newtheorem{exercise}[definition]{练习}
\hypersetup{unicode=true,colorlinks=true,linkcolor=blue}
\title{函数极限与连续：本科一年级微积分讲义计划}
\date{}
\begin{document}
\maketitle
\section{一元函数极限：定义、左右极限与无穷远极限}\label{sec:S01}
% UNIT S01.U01
\begin{definition}
设函数 $f$ 在点 $x_0$ 的某个去心邻域内有定义。若存在常数 $A$，使得对任意 $\varepsilon>0$，都存在 $\delta>0$，满足当 $0<|x-x_0|<\delta$ 时有 $|f(x)-A|<\varepsilon$，则称 $A$ 为 $f(x)$ 当 $x\to x_0$ 时的极限，记为
\[
\lim_{x\to x_0}f(x)=A.
\]
其中，$x_0$ 不要求属于函数的定义域。\label{def:point-limit}
\end{definition}
% UNIT S01.U02
\begin{definition}
设函数 $f$ 在 $x_0$ 的某个左侧去心邻域内有定义。若存在常数 $A_-$，使得对任意 $\varepsilon>0$，都存在 $\delta>0$，当 $x_0-\delta<x<x_0$ 时有 $|f(x)-A_-|<\varepsilon$，则称 $A_-$ 为 $f(x)$ 当 $x\to x_0^-$ 时的左极限，记为
\[
\lim_{x\to x_0^-}f(x)=A_-.
\]
设函数 $f$ 在 $x_0$ 的某个右侧去心邻域内有定义。若存在常数 $A_+$，使得对任意 $\varepsilon>0$，都存在 $\delta>0$，当 $x_0<x<x_0+\delta$ 时有 $|f(x)-A_+|<\varepsilon$，则称 $A_+$ 为 $f(x)$ 当 $x\to x_0^+$ 时的右极限，记为
\[
\lim_{x\to x_0^+}f(x)=A_+.
\]
\label{def:one-sided-limit}
\end{definition}
% UNIT S01.U03
\begin{definition}
设函数 $f$ 在某个右半无限区间 $(M,+\infty)$ 上有定义。若存在常数 $A$，使得对任意 $\varepsilon>0$，都存在 $R>M$，当 $x>R$ 时有 $|f(x)-A|<\varepsilon$，则称 $A$ 为 $f(x)$ 当 $x\to+\infty$ 时的极限，记为
\[
\lim_{x\to+\infty}f(x)=A.
\]
类似地，若函数在某个左半无限区间 $(-\infty,M)$ 上有定义，且对任意 $\varepsilon>0$ 都存在 $R<M$，使得当 $x<R$ 时有 $|f(x)-A|<\varepsilon$，则记为
\[
\lim_{x\to-\infty}f(x)=A.
\]
\label{def:infinite-limit}
\end{definition}
% UNIT S01.U04
\begin{theorem}
设函数 $f$ 在点 $x_0$ 的某个去心邻域内有定义，则
\[
\lim_{x\to x_0}f(x)=A
\]
当且仅当
\[
\lim_{x\to x_0^-}f(x)=A
\quad\text{且}\quad
\lim_{x\to x_0^+}f(x)=A.
\]
也就是说，函数在一点的双侧极限存在，当且仅当该点的左极限和右极限都存在且相等。\label{thm:two-sided-limit-criterion}
\end{theorem}
% UNIT S01.U05
\begin{proof}
先设 $\lim_{x\to x_0}f(x)=A$。给定任意 $\varepsilon>0$，由双侧极限的定义，存在 $\delta>0$，使得当 $0<|x-x_0|<\delta$ 时有 $|f(x)-A|<\varepsilon$。特别地，当 $x_0-\delta<x<x_0$ 或 $x_0<x<x_0+\delta$ 时，上述不等式分别成立，因此左极限和右极限都等于 $A$。

反过来，设左极限和右极限都等于 $A$。给定任意 $\varepsilon>0$，存在 $\delta_1>0$，使得当 $x_0-\delta_1<x<x_0$ 时有 $|f(x)-A|<\varepsilon$；存在 $\delta_2>0$，使得当 $x_0<x<x_0+\delta_2$ 时有 $|f(x)-A|<\varepsilon$。取 $\delta=\min\{\delta_1,\delta_2\}$，则当 $0<|x-x_0|<\delta$ 时，$x$ 必在 $x_0$ 的左侧或右侧，从而都有 $|f(x)-A|<\varepsilon$。因此 $\lim_{x\to x_0}f(x)=A$。
\end{proof}
% UNIT S01.U06
\begin{theorem}
设函数 $f$ 在点 $x_0$ 的某个去心邻域内有定义。若
\[
\lim_{x\to x_0}f(x)=A
\quad\text{且}\quad
\lim_{x\to x_0}f(x)=B,
\]
则 $A=B$。因此，函数在同一点处的有限极限至多有一个。\label{thm:limit-uniqueness}
\end{theorem}
% UNIT S01.U07
\begin{proof}
假设 $A\ne B$。取 $\varepsilon=|A-B|/3>0$。由 $\lim_{x\to x_0}f(x)=A$，存在 $\delta_1>0$，使得当 $0<|x-x_0|<\delta_1$ 时有 $|f(x)-A|<\varepsilon$；由 $\lim_{x\to x_0}f(x)=B$，存在 $\delta_2>0$，使得当 $0<|x-x_0|<\delta_2$ 时有 $|f(x)-B|<\varepsilon$。

取 $\delta=\min\{\delta_1,\delta_2\}$，并取一个满足 $0<|x-x_0|<\delta$ 的 $x$。由三角不等式，
\[
|A-B|\leq |A-f(x)|+|f(x)-B|<2\varepsilon=\frac{2}{3}|A-B|,
\]
这与 $|A-B|>0$ 矛盾。因此 $A=B$。
\end{proof}
% UNIT S01.U08
\begin{example}
计算
\[
\lim_{x\to2}(3x-1)\text{。}
\]
为说明双侧极限的处理方法，分别考察 $x\to2^-$ 和 $x\to2^+$ 两种情形。
\end{example}
% UNIT S01.U09
\par\noindent\textbf{解：}当 $x\to2^-$ 或 $x\to2^+$ 时，都有 $3x-1\to3\cdot2-1=5$。更直接地，对任意 $\varepsilon>0$，取 $\delta=\varepsilon/3$。当 $0<|x-2|<\delta$ 时，
\[
|(3x-1)-5|=|3x-6|=3|x-2|<3\delta=\varepsilon\text{。}
\]
因此由极限定义，
\[
\lim_{x\to2}(3x-1)=5\text{。}
\]\par
% UNIT S01.U10
\begin{example}
设
\[
f(x)=
\begin{cases}
1,&x<0,\\
2,&x>0.
\end{cases}
\]
判断 $\lim_{x\to0}f(x)$ 是否存在。
\end{example}
% UNIT S01.U11
\par\noindent\textbf{解：}当 $x\to0^-$ 时，始终有 $f(x)=1$，所以
\[
\lim_{x\to0^-}f(x)=1\text{。}
\]
当 $x\to0^+$ 时，始终有 $f(x)=2$，所以
\[
\lim_{x\to0^+}f(x)=2\text{。}
\]
由于左右极限不相等，根据定理 \ref{thm:two-sided-limit-criterion}，$\lim_{x\to0}f(x)$ 不存在。\par
% UNIT S01.U12
\begin{exercise}
\label{mathagent:exercise:bf60311c135cf0c14f3cb42069a931ed175758be7d89dff861bed83a8d332803}
设
\[
f(x)=
\begin{cases}
2x+1,&x<1,\\
5,&x>1.
\end{cases}
\]
判断 $\lim_{x\to1}f(x)$ 是否存在；若存在，求出该极限。
\end{exercise}
\section{极限的基本性质与运算}\label{sec:S02}
% UNIT S02.U01
\begin{definition}
设函数 $f,g,h$ 在点 $x_0$ 的某个去心邻域内有定义。若在该邻域内有 $h(x)\leq f(x)\leq g(x)$，且存在有限数 $L$，使得
\[
\lim_{x\to x_0}h(x)=L,\qquad \lim_{x\to x_0}g(x)=L,
\]
则
\[
\lim_{x\to x_0}f(x)=L\text{。}
\]
这称为夹逼原理。\label{def:squeeze-principle}
\end{definition}
% UNIT S02.U02
\begin{theorem}
设函数 $f$ 与 $g$ 在点 $x_0$ 的某个去心邻域内有定义，且 $\lim_{x\to x_0}f(x)=A$，$\lim_{x\to x_0}g(x)=B$。则
\[
\lim_{x\to x_0}(f(x)+g(x))=A+B,
\qquad
\lim_{x\to x_0}(f(x)-g(x))=A-B,
\]
\[
\lim_{x\to x_0}f(x)g(x)=AB.
\]
若 $B\neq0$，并且在某个去心邻域内 $g(x)\neq0$，则
\[
\lim_{x\to x_0}\frac{f(x)}{g(x)}=\frac{A}{B}.
\]
\label{thm:limit-algebra}

复合极限的基础结论：设 $g$ 在 $x_0$ 的某个去心邻域内有定义，$\lim_{x\to x_0}g(x)=b$；外函数 $f$ 在 $b$ 的某个去心邻域内有定义，且 $\lim_{t\to b}f(t)=L$。若在 $x_0$ 的某个去心邻域内，$g(x)\neq b$ 且 $g(x)$ 属于 $f$ 的定义域（因而复合函数有定义），则
\[
\lim_{x\to x_0}f(g(x))=L.
\]
这里要求内函数最终不取到 $b$，不能仅凭上述两个极限就无条件推出复合极限。
\end{theorem}
% UNIT S02.U03
\begin{proof}
先证明和的法则。给定 $\varepsilon>0$，由 $f(x)\to A$ 与 $g(x)\to B$，分别存在 $\delta_1>0$ 与 $\delta_2>0$，使得当 $0<|x-x_0|<\delta_1$ 时 $|f(x)-A|<\varepsilon/2$，当 $0<|x-x_0|<\delta_2$ 时 $|g(x)-B|<\varepsilon/2$。取 $\delta=\min\{\delta_1,\delta_2\}$，则
\[
\left|(f(x)+g(x))-(A+B)\right|\leq |f(x)-A|+|g(x)-B|<\varepsilon.
\]
故 $f(x)+g(x)\to A+B$。由于 $|(-g(x))-(-B)|=|g(x)-B|$，由定义有 $-g(x)\to-B$；差的法则由 $f(x)-g(x)=f(x)+(-g(x))$ 及和的法则得到。

先说明局部有界性：若 $f$ 在 $x_0$ 的某个去心邻域内有定义，且有有限极限 $A$，则 $f$ 在某个去心邻域内有界。事实上，在极限定义中取误差为 $1$，可得 $\delta_0>0$，使得当 $0<|x-x_0|<\delta_0$ 时，
\[
|f(x)-A|<1,\qquad |f(x)|\leq |A|+|f(x)-A|<|A|+1.
\]
这只涉及去心邻域，不要求 $f(x_0)$ 有定义。

对于积，令 $M=|A|+1>0$。给定 $\varepsilon>0$，由两个极限的定义可分别取 $\delta_1,\delta_2>0$，使得当 $0<|x-x_0|<\delta_1$ 时 $|g(x)-B|<\varepsilon/(2M)$，当 $0<|x-x_0|<\delta_2$ 时 $|f(x)-A|<\varepsilon/(2(|B|+1))$。取 $\delta=\min\{\delta_0,\delta_1,\delta_2\}$，则当 $0<|x-x_0|<\delta$ 时，
\[
|f(x)g(x)-AB|\leq |f(x)|\,|g(x)-B|+|B|\,|f(x)-A|
<\frac{\varepsilon}{2}+\frac{|B|\varepsilon}{2(|B|+1)}<\varepsilon.
\]
故由极限定义，$f(x)g(x)\to AB$。

若 $B\neq0$，由 $g(x)\to B$，可取 $\delta_3>0$，使得当 $0<|x-x_0|<\delta_3$ 时 $|g(x)-B|<|B|/2$。于是
\[
|g(x)|\geq |B|-|g(x)-B|>\frac{|B|}{2}>0,
\]
故分母不仅非零，而且远离零，并有
\[
\left|\frac{1}{g(x)}-\frac{1}{B}\right|=\frac{|g(x)-B|}{|g(x)||B|}\leq\frac{2|g(x)-B|}{|B|^2}.
\]
给定 $\varepsilon>0$，再由极限定义取 $\delta_4>0$，使得当 $0<|x-x_0|<\delta_4$ 时 $|g(x)-B|<\varepsilon|B|^2/2$。当 $0<|x-x_0|<\min\{\delta_3,\delta_4\}$ 时，上述倒数误差小于 $\varepsilon$。所以 $1/g(x)\to1/B$。再由已证明的积的法则，得到
\[
\frac{f(x)}{g(x)}=f(x)\frac{1}{g(x)}\to\frac{A}{B}.
\]

最后证明复合极限结论。按该结论的假设，取 $r>0$，使得当 $0<|x-x_0|<r$ 时，$g(x)\neq b$ 且 $f(g(x))$ 有定义。给定 $\varepsilon>0$，由 $\lim_{t\to b}f(t)=L$，存在 $\eta>0$，使得当 $0<|t-b|<\eta$ 时 $|f(t)-L|<\varepsilon$。由 $g(x)\to b$，存在 $\delta_5>0$，使得当 $0<|x-x_0|<\delta_5$ 时 $|g(x)-b|<\eta$。因此，当 $0<|x-x_0|<\min\{r,\delta_5\}$ 时，
\[
0<|g(x)-b|<\eta,\qquad |f(g(x))-L|<\varepsilon.
\]
由极限定义即得 $\lim_{x\to x_0}f(g(x))=L$。
\end{proof}
% UNIT S02.U04
\begin{theorem}
设函数 $f,g,h$ 在点 $x_0$ 的某个去心邻域内有定义，并且在该邻域内满足
\[
h(x)\leq f(x)\leq g(x).
\]
若
\[
\lim_{x\to x_0}h(x)=L,
\qquad
\lim_{x\to x_0}g(x)=L,
\]
则
\[
\lim_{x\to x_0}f(x)=L.
\]
\label{thm:squeeze-theorem}

夹逼计算的应用：求 $\lim_{x\to0}x^2\sin(1/x)$。构造上下界时，先利用三角函数的基本性质 $-1\leq\sin t\leq1$，再乘以非负因子 $x^2$。在去心邻域 $0<|x|<1$ 内，$1/x$ 有定义，且
\[
-x^2\leq x^2\sin(1/x)\leq x^2.
\]
两端的极限可以直接由定义核实：给定 $\varepsilon>0$，取 $\delta=\min\{1,\sqrt{\varepsilon}\}>0$。当 $0<|x|<\delta$ 时，
\[
|x^2-0|=x^2<\varepsilon,\qquad |(-x^2)-0|=x^2<\varepsilon.
\]
所以
\[
\lim_{x\to0}(-x^2)=0,\qquad \lim_{x\to0}x^2=0.
\]
因此由夹逼定理，
\[
\lim_{x\to0}x^2\sin(1/x)=0.
\]
这里无需计算 $\sin(1/x)$ 的极限，也未使用三角函数的重要极限；关键是用它的有界性构造趋于同一数值的上下界。
\end{theorem}
% UNIT S02.U05
\begin{proof}
给定 $\varepsilon>0$。由 $h(x)\to L$，存在 $\delta_1>0$，使得当 $0<|x-x_0|<\delta_1$ 时 $h(x)>L-\varepsilon$；由 $g(x)\to L$，存在 $\delta_2>0$，使得当 $0<|x-x_0|<\delta_2$ 时 $g(x)<L+\varepsilon$。取 $\delta=\min\{\delta_1,\delta_2\}$，则当 $0<|x-x_0|<\delta$ 时，
\[
L-\varepsilon<h(x)\leq f(x)\leq g(x)<L+\varepsilon.
\]
因此 $|f(x)-L|<\varepsilon$，从而 $\lim_{x\to x_0}f(x)=L$。
\end{proof}
% UNIT S02.U06
\begin{theorem}
设函数 $f$ 在点 $x_0$ 的某个去心邻域内有定义，且 $\lim_{x\to x_0}f(x)=A$。若 $A>0$，则存在 $\delta>0$，使得当 $0<|x-x_0|<\delta$ 时 $f(x)>0$；若 $A<0$，则存在 $\delta>0$，使得当 $0<|x-x_0|<\delta$ 时 $f(x)<0$。\label{thm:limit-sign-preservation}
\end{theorem}
% UNIT S02.U07
\begin{proof}
设 $A>0$。取 $\varepsilon=A/2$。由极限定义，存在 $\delta>0$，使得当 $0<|x-x_0|<\delta$ 时 $|f(x)-A|<A/2$，于是
\[
f(x)>A-\frac{A}{2}=\frac{A}{2}>0.
\]
当 $A<0$ 时，取 $\varepsilon=|A|/2$，则
\[
f(x)<A+\frac{|A|}{2}=\frac{A}{2}<0.
\]

\end{proof}
% UNIT S02.U08
\begin{example}
计算
\[
\lim_{x\to1}\frac{x^2-1}{x-1}.
\]
要求在保持定义域条件 $x\neq1$ 的前提下进行代数变形。
\end{example}
% UNIT S02.U09
\par\noindent\textbf{解：}当 $x\neq1$ 时，$x^2-1=(x-1)(x+1)$，所以
\[
\frac{x^2-1}{x-1}=x+1.
\]
原函数虽然在 $x=1$ 处没有定义，但极限只考察 $x\neq1$ 且 $x\to1$ 时的函数值。这里直接用定义核实约分后函数的极限：给定任意 $\varepsilon>0$，取 $\delta=\varepsilon>0$。当 $0<|x-1|<\delta$ 时，
\[
|(x+1)-2|=|x-1|<\varepsilon.
\]
故 $\lim_{x\to1}(x+1)=2$。在同一去心邻域内，原函数与 $x+1$ 相等，其与 $2$ 的误差也小于 $\varepsilon$，因此
\[
\lim_{x\to1}\frac{x^2-1}{x-1}=\lim_{x\to1}(x+1)=2.
\]\par
% UNIT S02.U10
\begin{example}
设 $f(x)=x^2-3x+2$。求 $\lim_{x\to2}f(x)$，并说明所得极限的符号。
\end{example}
% UNIT S02.U11
\par\noindent\textbf{解：}先由极限定义核实运算法则所需的基本极限。给定 $\varepsilon>0$，取 $\delta=\varepsilon$，则当 $0<|x-2|<\delta$ 时，恒等函数与目标值 $2$ 的误差为 $|x-2|<\varepsilon$，故 $\lim_{x\to2}x=2$。对常数函数 $c$，取任意 $\delta>0$ 都有 $|c-c|=0<\varepsilon$，故 $\lim_{x\to2}c=c$；特别地，常数函数 $3$ 和 $2$ 的极限分别为 $3$ 和 $2$。

由上述基本极限及极限的和、差与积法则，
\[
\lim_{x\to2}(x^2-3x+2)=2^2-3\cdot2+2=0.
\]
因此该极限等于 $0$。由于极限值为零，保号定理不能据此断言 $f(x)$ 在 $2$ 的去心邻域内恒正或恒负。\par
% UNIT S02.U12
\begin{exercise}
\label{mathagent:exercise:b4c9c9a1fb82bcef052fdbf6e1539831ed88cee016dbcbf83a50e5e8d2311be0}
设函数 $f$ 与 $g$ 在点 $x_0$ 的某个去心邻域内有定义，且
\[
\lim_{x\to x_0}f(x)=2,
\qquad
\lim_{x\to x_0}g(x)=-1.
\]
求
\[
\lim_{x\to x_0}\frac{f(x)+g(x)}{f(x)-g(x)}.
\]
并说明使用商的极限法则时需要检查的条件。

附加小问：用夹逼定理计算
\[
\lim_{x\to0}\frac{x^2}{2+\cos(1/x)}.
\]
请自行构造上下界，确认原式与上下界在同一去心邻域内有定义，说明所考察的同一趋近过程，计算两端的共同极限，再给出结论。
\end{exercise}
\section{无穷小、基本极限与等价无穷小}\label{sec:S03}
% UNIT S03.U01
\begin{definition}
设函数 $\alpha(x)$ 在考察点 $x_0$ 的某个去心邻域内有定义。若
\[
\lim_{x\to x_0}\alpha(x)=0,
\]
则称 $\alpha(x)$ 是当 $x\to x_0$ 时的无穷小。这里的“无穷小”表示函数值趋近于零，而不表示函数在某个邻域内恒不为零。若对任意 $M>0$，都存在 $\delta>0$，使得当 $0<|x-x_0|<\delta$ 时有
\[
|\alpha(x)|>M,
\]
则称 $\alpha(x)$ 是当 $x\to x_0$ 时绝对值意义下的无穷大，即 $|\alpha(x)|\to+\infty$。本节的“无穷大”采用这一含义，它本身不指定函数值的正负方向。

正无穷极限与负无穷极限须分别定义：若对任意 $M>0$，都存在 $\delta>0$，使得当 $0<|x-x_0|<\delta$ 时有 $\alpha(x)>M$，则记为
\[
\lim_{x\to x_0}\alpha(x)=+\infty;
\]
若对任意 $M>0$，都存在 $\delta>0$，使得当 $0<|x-x_0|<\delta$ 时有 $\alpha(x)<-M$，则记为
\[
\lim_{x\to x_0}\alpha(x)=-\infty.
\]
上述两种情形均满足绝对值意义下的无穷大条件，但仅有 $|\alpha(x)|\to+\infty$ 不能推出其中任何一种带符号的极限。无穷小与无穷大都是相对于特定的极限过程而言的。\label{def:infinitesimal}
\end{definition}
% UNIT S03.U02
\begin{definition}
设函数 $\alpha(x)$ 与 $\beta(x)$ 在 $x_0$ 的某个去心邻域内有定义，且在该邻域内 $\beta(x)\neq0$。若
\[
\lim_{x\to x_0}\alpha(x)=0,\qquad
\lim_{x\to x_0}\beta(x)=0,
\qquad
\lim_{x\to x_0}\frac{\alpha(x)}{\beta(x)}=1,
\]
则称 $\alpha(x)$ 与 $\beta(x)$ 当 $x\to x_0$ 时为等价无穷小，记为
\[
\alpha(x)\sim\beta(x)\qquad(x\to x_0).
\]
其中，前两个条件说明二者都是同一极限过程下的无穷小，最后一个条件给出它们的等价关系；比值条件只在商有定义的去心邻域内使用。若只满足比值极限为 $1$ 而不要求二者都趋于零，则称为一般的等价函数关系，而不称为等价无穷小。\label{def:equivalent-infinitesimal}

无穷小的阶比较描述它们在同一极限过程中趋于零的相对快慢。仍设 $\alpha(x)\to0$、$\beta(x)\to0$，且 $\beta(x)$ 在某个去心邻域内非零。当 $x\to x_0$ 时，若
\[
\frac{\alpha(x)}{\beta(x)}\to0,
\]
则称 $\alpha$ 是比 $\beta$ 高阶的无穷小，记为 $\alpha=o(\beta)$；若
\[
\left|\frac{\alpha(x)}{\beta(x)}\right|\to+\infty,
\]
则称 $\alpha$ 是比 $\beta$ 低阶的无穷小；若
\[
\frac{\alpha(x)}{\beta(x)}\to c,\qquad 0<|c|<\infty,
\]
则称二者为同阶无穷小。因此，等价无穷小是同阶无穷小中 $c=1$ 的特殊情形，并非所有同阶无穷小都等价。这些比值条件不必对任意一对无穷小成立，不能在比值没有相应极限时强行作上述分类。

记号 $r(x)=o(1)$ 表示在指定的极限过程中 $r(x)\to0$；这里的 $1$ 是比较用的常数，并不是无穷小。在上述分母非零的条件下，$\alpha=o(\beta)$ 等价于 $\alpha=\beta r$，其中 $r=o(1)$；$\alpha\sim\beta$ 则等价于 $\alpha=\beta(1+r)$，其中 $r=o(1)$。这些记号均须注明或由上下文明确同一个极限过程。
\end{definition}
% UNIT S03.U03
\begin{theorem}
在实数范围内，当 $x\to0$ 时，有两个基本极限
\[
\lim_{x\to0}\frac{\sin x}{x}=1,
\qquad
\lim_{x\to0}\frac{e^x-1}{x}=1.
\]
其中，第一式中的 $x$ 取弧度制。由此可得
\[
\sin x\sim x,
\qquad
 e^x-1\sim x\qquad(x\to0).
\]
这些结论只适用于所写出的极限过程及其定义域条件，不能不加条件地推广到任意点或任意角度单位。\label{thm:fundamental-limits}
\end{theorem}
% UNIT S03.U04
\begin{proof}
先补充说明三角函数部分所需的基础事实。这里所有角均采用弧度制。由单位圆中的几何不等式，对任意实数 $u$ 都有
\[
|\sin u|\leq |u|.
\]
因此取 $u=x/2$，并结合三角恒等式，得到
\[
1-\cos x=2\sin^2\frac{x}{2},
\qquad
0\leq1-\cos x=2\sin^2\frac{x}{2}\leq2\left|\frac{x}{2}\right|^2=\frac{x^2}{2}.
\]
这里的估计在 $x\to0$ 的某个邻域内当然成立，实际上对所有实数 $x$ 都成立。由于 $x^2/2\to0$，由夹逼原理得到
\[
\lim_{x\to0}\cos x=1.
\]
因此，下面使用的 $\cos x\to1$ 已由上述估计得到，而不是未经说明的前提。

对 $0<x<\dfrac{\pi}{2}$，由单位圆中的几何不等式，有
\[
\sin x<x<\tan x.
\]
由于 $\sin x>0$，两边分别除以 $\sin x$，得到
\[
1<\frac{x}{\sin x}<\frac{1}{\cos x}.
\]
取倒数可得
\[
\cos x<\frac{\sin x}{x}<1.
\]
当 $x\to0^+$ 时，$\cos x\to1$。这里所用的是夹逼原理在右侧邻域上的同样结论：若在某个右侧去心邻域内有 $h(x)\leq f(x)\leq g(x)$，且当 $x\to x_0^+$ 时 $h(x)$ 与 $g(x)$ 都趋于同一有限数 $L$，则对任意 $\varepsilon>0$，分别由两个单侧极限取得适用的 $\delta_1,\delta_2$，取 $\delta=\min\{\delta_1,\delta_2\}$，即可得到 $|f(x)-L|<\varepsilon$；因此 $f(x)\to L$。这只是双侧夹逼证明的同样 $\delta$ 论证，并将考察范围限制在右侧邻域。由此可知
\[
\lim_{x\to0^+}\frac{\sin x}{x}=1.
\]
又因为 $\sin x/x$ 是偶函数，所以左极限也等于 $1$，从而
\[
\lim_{x\to0}\frac{\sin x}{x}=1.
\]

再证明第二个基本极限。先建立所需的指数函数估计，而不把待证极限作为先验性质。对正整数 $n$，设
\[
a_n=\left(1+\frac1n\right)^n.
\]
由二项式定理，$a_n$ 的第 $k$ 项为
\[
\binom nk\frac1{n^k}=\frac1{k!}\prod_{j=0}^{k-1}\left(1-\frac{j}{n}\right).
\]
对每个固定的 $k$，该项随 $n$ 增大而不减；展开式中的新增项非负，所以 $a_n$ 单调不减。又因 $k!\geq2^{k-1}$（$k\geq1$），当 $n\geq2$ 时
\[
a_n\leq1+1+\sum_{k=2}^n\frac1{2^{k-1}}<3,
\]
且 $a_1=2$。令
\[
E=\sup\{a_n:n\in\mathbb N,\ n\geq1\}.
\]
这里使用实数的上确界性质：每个非空且有上界的实数集合都有最小上界。由上面的估计，集合 $\{a_n:n\geq1\}$ 非空且有上界 $3$，故 $E$ 存在，并且 $2\leq E\leq3$。

现在用数列极限的 $\varepsilon$-$N$ 含义说明 $a_n$ 确实趋于 $E$。任给 $\varepsilon>0$，由 $E$ 是上确界，数 $E-\varepsilon$ 不可能仍是该集合的上界，因此存在正整数 $N$ 使得
\[
a_N>E-\varepsilon.
\]
由于 $(a_n)$ 单调不减且 $a_n\leq E$，对一切 $n\geq N$ 都有
\[
0\leq E-a_n\leq E-a_N<\varepsilon.
\]
这正是 $a_n\to E$ 的 $\varepsilon$-$N$ 定义。因此可定义
\[
e:=E=\lim_{n\to\infty}a_n,
\qquad e\geq2>1.
\]
同一 $\varepsilon$-$N$ 论证还说明任意子列 $a_{qn}$（其中 $q$ 为固定正整数）也趋于 $e$：给定 $\varepsilon>0$，取上面得到的 $N$，当 $n\geq N$ 时有 $qn\geq N$，从而
\[
|a_{qn}-e|<\varepsilon.
\]

实数幂先由有理幂的保序延拓定义。下面说明所需的根式极限，而不把它作为未说明的连续性结论使用。固定正整数 $q$，设 $y_n>0$、$y>0$ 且 $y_n\to y$。令 $r_n=y_n^{1/q}$、$r=y^{1/q}$，则
\[
y_n-y=(r_n-r)(r_n^{q-1}+r_n^{q-2}r+\cdots+r^{q-1}).
\]
由于 $y_n\to y>0$，当 $n$ 足够大时 $y_n>y/2$，故 $r_n\geq(y/2)^{1/q}$；因而上式中的正因子之和有一个与 $n$ 无关的正下界。这说明
\[
|r_n-r|\leq C|y_n-y|\to0
\]
（其中 $C$ 为某个正常数），所以 $y_n^{1/q}\to y^{1/q}$。再由有限因式分解，对固定正整数 $p$ 有
\[
r_n^p-r^p=(r_n-r)(r_n^{p-1}+r_n^{p-2}r+\cdots+r^{p-1}),
\]
从而 $y_n^{p/q}=r_n^p\to r^p=y^{p/q}$。这就明确给出了正数有理幂所需的极限依据。

对有理数 $t=p/q\in(0,1)$（$p,q$ 为正整数），取子列 $n\mapsto qn$，有
\[
a_{qn}^t=\left(1+\frac1{qn}\right)^{pn}.
\]
由上面已经证明的子列极限和根式误差控制，$a_{qn}^t\to e^t$。二项式展开给出下界
\[
a_{qn}^t\geq1+\frac{pn}{qn}=1+t.
\]
另一方面，对每个 $k$，有
\[
\frac{\binom{pn}{k}}{(qn)^k}
\leq\frac{(pn)^k}{k!(qn)^k}
=\frac{t^k}{k!}\leq t^k,
\]
从而
\[
a_{qn}^t\leq\sum_{k=0}^{pn}t^k
=\frac{1-t^{pn+1}}{1-t}
<\frac1{1-t}.
\]
令 $n\to\infty$，得到对有理 $t\in(0,1)$ 有
\[
1+t\leq e^t\leq\frac1{1-t}.
\]
这里使用的是有限等比求和公式，没有另行使用无穷级数收敛理论。

对任意实数 $0<t<1$，取有理数 $0<u<t<v<1$。由实幂的保序性，
\[
1+u\leq e^u\leq e^t\leq e^v\leq\frac1{1-v}.
\]
令有理数 $u$ 从下方趋于 $t$、$v$ 从上方趋于 $t$。这里只对线性函数和有理式取极限，得到
\[
1+t\leq e^t\leq\frac1{1-t},\qquad 0<t<1.
\]
这就建立了后续证明所需的指数函数估计，未使用待证的差商极限、指数函数在零点的导数或 Taylor 展开。

当 $0<x<1$ 时，由上述不等式得到
\[
1\leq\frac{e^x-1}{x}\leq\frac{1}{1-x}.
\]
令 $x\to0^+$，并应用上面说明的右侧夹逼结论，得到
\[
\lim_{x\to0^+}\frac{e^x-1}{x}=1.
\]
当 $x=-t<0$ 且 $0<t<1$ 时，对
\[
1+t\leq e^t\leq\frac1{1-t}
\]
取倒数，得到
\[
1-t\leq e^{-t}\leq\frac1{1+t}.
\]
由于 $-t<0$，相减并除以负数 $-t$ 时不等号方向反转，因而
\[
\frac1{1+t}\leq\frac{e^{-t}-1}{-t}\leq1.
\]
令 $t\to0^+$，再次应用右侧夹逼结论，得到
\[
\lim_{x\to0^-}\frac{e^x-1}{x}=1.
\]
左右极限相等，所以
\[
\lim_{x\to0}\frac{e^x-1}{x}=1.
\]
整个论证只使用已说明的三角函数估计、实数上确界性质、数列极限的 $\varepsilon$-$N$ 含义、由 Euler 定义和有理幂根式估计建立的指数函数不等式以及夹逼定理。
\end{proof}
% UNIT S03.U05
\begin{theorem}
设 $\alpha(x)$ 与 $\beta(x)$ 在 $x_0$ 的某个去心邻域内有定义，且在该邻域内均不为零。若
\[
\alpha(x)\sim\beta(x)\qquad(x\to x_0),
\]
并且函数 $F(x)$ 在相应去心邻域内有定义，满足 $F(x)\alpha(x)$ 有定义以及
\[
\lim_{x\to x_0}F(x)\alpha(x)=L,
\]
其中 $L$ 为有限数，则
\[
\lim_{x\to x_0}F(x)\beta(x)=L,
\]
只要 $F(x)\beta(x)$ 在该去心邻域内有定义。特别地，在乘积的极限中，可以用与某个无穷小等价的无穷小替换它，但必须保持相关表达式有定义。

商情形也有明确的替换条件：若 $F(x)/\alpha(x)$ 在某个去心邻域内有定义，且
\[
\lim_{x\to x_0}\frac{F(x)}{\alpha(x)}=L,
\]
其中 $L$ 为有限数，并且 $F(x)/\beta(x)$ 在同一去心邻域内有定义，则
\[
\lim_{x\to x_0}\frac{F(x)}{\beta(x)}=L.
\]
这里要求 $\alpha(x)$ 与 $\beta(x)$ 在该去心邻域内均不为零，从而 $\alpha(x)/\beta(x)$ 有定义；证明为
\[
\frac{F(x)}{\beta(x)}=\frac{F(x)}{\alpha(x)}\frac{\alpha(x)}{\beta(x)}.
\]
由 $\alpha(x)\sim\beta(x)$ 可知
\[
\lim_{x\to x_0}\frac{\alpha(x)}{\beta(x)}=1.
\]
再由极限的乘法法则，得到
\[
\lim_{x\to x_0}\frac{F(x)}{\beta(x)}
=\left(\lim_{x\to x_0}\frac{F(x)}{\alpha(x)}\right)
\left(\lim_{x\to x_0}\frac{\alpha(x)}{\beta(x)}\right)
=L\cdot1=L.
\]
因此，在商的极限中，可以将分母中的无穷小 $\alpha(x)$ 替换为与其等价的无穷小 $\beta(x)$，但必须同时检查两个商在同一去心邻域内有定义，并且原商的有限极限存在。一般而言，不能在加减式中无条件逐项替换等价无穷小；加减项的替换必须另行检验整体极限，不能仅由等价关系和表达式有定义推出。\label{thm:equivalent-substitution}
\end{theorem}
% UNIT S03.U06
\begin{proof}
由 $\alpha(x)\sim\beta(x)$ 的定义，有
\[
\lim_{x\to x_0}\frac{\alpha(x)}{\beta(x)}=1.
\]
又由等价无穷小替换定理的假设，$\alpha(x)$ 与 $\beta(x)$ 在同一去心邻域内均不为零，因此商 $\beta(x)/\alpha(x)$ 在该邻域内有定义，并且
\[
\frac{\beta(x)}{\alpha(x)}
=\left(\frac{\alpha(x)}{\beta(x)}\right)^{-1}.
\]
由于 $\lim_{x\to x_0}\alpha(x)/\beta(x)=1\neq0$，由倒数极限定理可得
\[
\lim_{x\to x_0}\frac{\beta(x)}{\alpha(x)}=1.
\]
在所给去心邻域内，利用乘积的代数分解，
\[
F(x)\beta(x)=F(x)\alpha(x)\frac{\beta(x)}{\alpha(x)}.
\]
由假设和极限的乘法法则，得到
\[
\lim_{x\to x_0}F(x)\beta(x)
=\left(\lim_{x\to x_0}F(x)\alpha(x)\right)
\left(\lim_{x\to x_0}\frac{\beta(x)}{\alpha(x)}\right)
=L\cdot1=L.
\]
因此，在上述条件下，若原式写成 $F(x)\alpha(x)$，可以把其中的因子 $\alpha(x)$ 替换为 $\beta(x)$，得到 $F(x)\beta(x)$。这一操作等价于将原式整体乘以修正比值 $\beta(x)/\alpha(x)$，而不是用该比值替换原式中的某个因子。替换时，$\alpha(x)$ 与 $\beta(x)$ 在同一去心邻域内均须非零，相关表达式须有定义，并且 $\lim_{x\to x_0}F(x)\alpha(x)$ 必须存在且为有限数。该结论适用于所证明的乘积结构，以及通过倒数改写后满足同样条件的商结构。它不表示可以在加减式中无条件逐项替换等价无穷小；例如加减项的替换必须另行检验整体极限，不能仅由等价关系和表达式有定义推出。
\end{proof}
% UNIT S03.U07
\begin{example}
计算
\[
\lim_{x\to0}\frac{(1+x)\sin x}{x}.
\]
要求说明等价无穷小替换的依据及适用条件。

另考察加减式中逐项替换的风险：当 $x\to0$ 且 $x\neq0$ 时，求
\[
\lim_{x\to0}\frac{(x+x^2)-x}{x^2},
\]
并与把加法项 $x+x^2$ 错误替换为等价无穷小 $x$ 后所得的表达式进行比较。
\end{example}
% UNIT S03.U08
\par\noindent\textbf{解：}当 $x\to0$ 时，$\sin x\sim x$，并且 $1+x\to1$。在 $x\neq0$ 的去心邻域内，原式有定义，且
\[
\frac{(1+x)\sin x}{x}=(1+x)\frac{\sin x}{x}.
\]
由基本极限
\[
\lim_{x\to0}\frac{\sin x}{x}=1,
\qquad
\lim_{x\to0}(1+x)=1,
\]
可得
\[
\lim_{x\to0}\frac{(1+x)\sin x}{x}=1\cdot1=1.
\]
这里的替换只在 $x\to0$ 且 $x\neq0$ 的去心邻域内进行，未将原式在 $x=0$ 处补定义。

再说明加减式中不能无条件逐项替换等价无穷小。对 $x\to0$ 且 $x\neq0$，有
\[
\frac{x+x^2}{x}=1+x\longrightarrow1,
\]
所以 $x+x^2\sim x$；同时分母 $x^2\neq0$。但是
\[
\frac{(x+x^2)-x}{x^2}=\frac{x^2}{x^2}=1,
\]
故原式的极限为 $1$。若把加减项 $x+x^2$ 逐项替换为等价无穷小 $x$，则得到
\[
\frac{x-x}{x^2}=0,
\]
其极限为 $0$，与原式不同。这里被替换的是加减式中的一个加项，而不是等价无穷小替换定理所许可的整体乘积因子或商结构；因此不能仅由等价关系和表达式有定义就进行这种替换，必须另行检验替换后的整体极限。\par
% UNIT S03.U09
\begin{exercise}
\label{mathagent:exercise:410f498083ff9a28d1a2acf9839e74554b434d3445561642fdfe5d3854f06971}
计算
\[
\lim_{x\to0}\frac{e^x-1}{x}\,(1-2x),
\]
并说明所使用的基本极限及乘法法则的条件。
\end{exercise}
\section{函数连续性与间断点}\label{sec:S04}
% UNIT S04.U01
\begin{definition}
设函数 $f$ 在点 $x_0$ 的某个邻域内有定义。若 $f$ 在 $x_0$ 处有定义，且
\[
\lim_{x\to x_0}f(x)=f(x_0)\text{，}
\]
则称函数 $f$ 在点 $x_0$ 处连续。等价地，$f$ 在 $x_0$ 处连续是指：对任意 $\varepsilon>0$，都存在 $\delta>0$，使得当 $|x-x_0|<\delta$ 时有
\[
|f(x)-f(x_0)|<\varepsilon.
\]
其中，条件 $f(x_0)$ 有定义不可省略。若 $f$ 在区间内每一点都连续，则称 $f$ 在该区间上连续。\label{def:point-continuity}
\end{definition}
% UNIT S04.U02
\begin{definition}
设 $f$ 在区间 $I$ 上有定义。若对每个内点 $x_0\in I$，函数 $f$ 在 $x_0$ 处连续；当 $I$ 含有端点时，函数在左端点处连续是指右极限等于函数值，在右端点处连续是指左极限等于函数值，则称 $f$ 在区间 $I$ 上连续。

特别地，当 $I=[a,b]$ 且 $a<b$ 时，$f$ 在 $[a,b]$ 上连续意味着：对每个 $x_0\in(a,b)$，有
\[
\lim_{x\to x_0}f(x)=f(x_0)\text{，}
\]
并且有
\[
\lim_{x\to a^+}f(x)=f(a)\text{，}
\qquad
\lim_{x\to b^-}f(x)=f(b)\text{。}
\]
\label{def:interval-continuity}
\end{definition}
% UNIT S04.U03
\begin{definition}
设函数 $f$ 在点 $x_0$ 的某个双侧去心邻域内有定义。若 $f(x_0)$ 未定义，或 $f(x_0)$ 有定义但 $f$ 在 $x_0$ 处不连续，则称 $x_0$ 为 $f$ 的间断点。因此，间断点不要求属于函数的定义域。下述间断点分类均保留这一双侧去心邻域条件，单侧无穷极限的独立定义则不受此条件限制。常见类型如下。

若极限 $\lim_{x\to x_0}f(x)$ 存在且为有限数，但 $f(x_0)$ 未定义或 $f(x_0)\neq\lim_{x\to x_0}f(x)$，则称 $x_0$ 为可去间断点。

若左、右极限都存在且有限，但不相等，即
\[
\lim_{x\to x_0^-}f(x)\neq\lim_{x\to x_0^+}f(x)\text{，}
\]
则称 $x_0$ 为跳跃间断点。

为说明无穷间断点，先独立定义单侧无穷极限；这些定义不要求上述双侧去心邻域条件。定义右侧无穷极限时，只设 $f$ 在某个右侧去心邻域 $(x_0,x_0+r)$ 内有定义，其中 $r>0$。若对任意 $M>0$，都存在 $\delta>0$，使得当 $0<x-x_0<\delta$ 时有 $f(x)>M$，则记为
\[
\lim_{x\to x_0^+}f(x)=+\infty.
\]
若对任意 $M>0$，都存在 $\delta>0$，使得当 $0<x-x_0<\delta$ 时有 $f(x)<-M$，则记为
\[
\lim_{x\to x_0^+}f(x)=-\infty.
\]
定义左侧无穷极限时，只设 $f$ 在某个左侧去心邻域 $(x_0-r,x_0)$ 内有定义，其中 $r>0$，不要求右侧有定义。若对任意 $M>0$，都存在 $\delta>0$，使得当 $0<x_0-x<\delta$ 时有 $f(x)>M$，则记为
\[
\lim_{x\to x_0^-}f(x)=+\infty.
\]
若对任意 $M>0$，都存在 $\delta>0$，使得当 $0<x_0-x<\delta$ 时有 $f(x)<-M$，则记为
\[
\lim_{x\to x_0^-}f(x)=-\infty.
\]
这里的 $\delta$ 均可取小到使相应单侧去心邻域包含在函数的定义域内；右侧定义不要求左侧有定义，左侧定义不要求右侧有定义，且均不要求 $f(x_0)$ 有定义。这些单侧极限定义本身并不把仅在单侧有定义的边界点归为本节所称的间断点。

例如，考虑定义在 $\mathbb{R}\setminus\{0\}$ 上的函数
\[
f(x)=\frac{1}{x}.
\]
在 $x_0=0$ 处考察右侧。对任意 $M>0$，取 $\delta=1/M$，当 $0<x<\delta$ 时有 $f(x)=1/x>M$，所以
\[
\lim_{x\to0^+}\frac{1}{x}=+\infty.
\]
由于该函数在 $0$ 的某个双侧去心邻域内有定义，且右极限为 $+\infty$，因此，$x=0$ 是无穷间断点。这里明确使用了定义域 $\mathbb{R}\setminus\{0\}$、考察方向 $x\to0^+$ 以及“对任意 $M>0$，存在 $\delta>0$”的量词条件。

在函数 $f$ 于 $x_0$ 的某个双侧去心邻域内有定义的前提下，若至少有一个单侧极限为 $+\infty$ 或 $-\infty$，即函数值沿相应方向趋于正无穷或负无穷，则称 $x_0$ 为无穷间断点。本讲义将可去间断点与跳跃间断点合称为第一类间断点，将其余间断点称为第二类间断点。等价地，第二类间断点是至少有一个单侧极限不存在为有限数的间断点；因而无穷间断点属于第二类间断点，但并非第二类间断点都为无穷间断点。若一个间断点既不属于可去或跳跃类型，又没有任何一个单侧极限为 $+\infty$ 或 $-\infty$，则它属于第二类间断点，但不是这里定义的无穷间断点。仅凭函数在每个双侧去心邻域内无界，不能断定它满足某个单侧无穷极限的条件，本节也不据此赋予它“无界振荡”的名称。分类时必须先确认函数在某个双侧去心邻域内有定义，再核对所考察的极限存在性条件；仅在单侧邻域有定义时，可以讨论相应的单侧无穷极限，但不纳入本节上述间断点分类。\label{def:discontinuity-types}
\end{definition}
% UNIT S04.U04
\begin{theorem}
设函数 $f$ 与 $g$ 在点 $x_0$ 处连续，且相关运算在 $x_0$ 处有定义。则函数
\[
f+g,
\qquad f-g,
\qquad fg
\]
在 $x_0$ 处连续。若 $g(x_0)\neq0$，并且在 $x_0$ 的某个邻域内 $g(x)\neq0$，则
\[
\frac{f}{g}
\]
在 $x_0$ 处连续。

若函数 $g$ 在 $x_0$ 处连续，函数 $f$ 在 $g(x_0)$ 处连续，并且复合函数 $f\circ g$ 在 $x_0$ 的某个邻域内有定义，则 $f\circ g$ 在 $x_0$ 处连续。\label{thm:continuity-algebra}
\end{theorem}
% UNIT S04.U05
\begin{proof}
由连续性的定义，函数 $f$ 与 $g$ 在 $x_0$ 处连续分别意味着
\[
\lim_{x\to x_0}f(x)=f(x_0)\text{，}
\qquad
\lim_{x\to x_0}g(x)=g(x_0)\text{。}
\]
根据极限的和、差与积法则，有
\[
\lim_{x\to x_0}(f(x)+g(x))=f(x_0)+g(x_0)=(f+g)(x_0)\text{，}
\]
\[
\lim_{x\to x_0}(f(x)-g(x))=f(x_0)-g(x_0)=(f-g)(x_0)\text{，}
\]
以及
\[
\lim_{x\to x_0}f(x)g(x)=f(x_0)g(x_0)=(fg)(x_0)\text{。}
\]
因此和、差与积在 $x_0$ 处连续。

若 $g(x_0)\neq0$ 且 $g$ 在 $x_0$ 的某个邻域内不为零，则由商的极限法则，
\[
\lim_{x\to x_0}\frac{f(x)}{g(x)}
=\frac{f(x_0)}{g(x_0)}
=\left(\frac{f}{g}\right)(x_0)\text{，}
\]
所以商在 $x_0$ 处连续。

最后设 $g$ 在 $x_0$ 处连续，$f$ 在 $g(x_0)$ 处连续，并且存在 $x_0$ 的某个邻域 $U$，使得对每个 $x\in U$，$g(x)$ 都属于 $f$ 的定义域。令 $\varepsilon>0$。由于 $f$ 在 $g(x_0)$ 处连续，存在 $\eta>0$，使得当 $|t-g(x_0)|<\eta$ 时，
\[
|f(t)-f(g(x_0))|<\varepsilon.
\]
由于 $g$ 在 $x_0$ 处连续，存在 $\delta_1>0$，使得当 $|x-x_0|<\delta_1$ 时，
\[
|g(x)-g(x_0)|<\eta.
\]
又因为 $U$ 是 $x_0$ 的邻域，存在 $\delta_2>0$，使得当 $|x-x_0|<\delta_2$ 时，$x\in U$。取 $\delta=\min\{\delta_1,\delta_2\}$。当 $|x-x_0|<\delta$ 时，$g(x)$ 属于 $f$ 的定义域，且
\[
|g(x)-g(x_0)|<\eta.
\]
令 $t=g(x)$，便有
\[
|f(g(x))-f(g(x_0))|<\varepsilon.
\]
由于 $(f\circ g)(x)=f(g(x))$ 且 $(f\circ g)(x_0)=f(g(x_0))$，故
\[
|(f\circ g)(x)-(f\circ g)(x_0)|<\varepsilon.
\]
由点连续的定义，$f\circ g$ 在 $x_0$ 处连续。
\end{proof}
% UNIT S04.U06
\begin{theorem}
设函数 $g$ 在点 $x_0$ 处连续，函数 $f$ 在点 $g(x_0)$ 处连续，并且存在 $x_0$ 的某个邻域 $U$，使得对每个 $x\in U$，$g(x)$ 都属于 $f$ 的定义域。则复合函数
\[
(f\circ g)(x)=f(g(x))
\]
在点 $x_0$ 处连续。\label{thm:composition-continuity}
\end{theorem}
% UNIT S04.U07
\begin{proof}
令 $\varepsilon>0$。由于 $f$ 在 $g(x_0)$ 处连续，存在 $\eta>0$，使得当 $|t-g(x_0)|<\eta$ 时，
\[
|f(t)-f(g(x_0))|<\varepsilon.
\]
由于 $g$ 在 $x_0$ 处连续，存在 $\delta>0$，使得当 $|x-x_0|<\delta$ 时，
\[
|g(x)-g(x_0)|<\eta.
\]
根据假设，当 $x$ 属于该邻域 $U$ 时，$g(x)$ 属于 $f$ 的定义域；缩小 $\delta$ 后还可令 $|x-x_0|<\delta$ 的点属于 $U$。于是对满足 $|x-x_0|<\delta$ 的 $x$，令 $t=g(x)$，便有
\[
|t-g(x_0)|=|g(x)-g(x_0)|<\eta
\]
以及
\[
|f(g(x))-f(g(x_0))|<\varepsilon.
\]
而
\[
(f\circ g)(x_0)=f(g(x_0))\text{，}
\]
故
\[
|(f\circ g)(x)-(f\circ g)(x_0)|<\varepsilon.
\]
由点连续的定义，$f\circ g$ 在 $x_0$ 处连续。
\end{proof}
% UNIT S04.U08
\begin{example}
判断函数
\[
f(x)=\frac{x^2-1}{x-1}
\]
在 $x=1$ 处的连续性，并指出该点的间断类型。
\end{example}
% UNIT S04.U09
\par\noindent\textbf{解：}函数 $f$ 的定义域为 $\mathbb{R}\setminus\{1\}$，因此 $f(1)$ 未定义。当 $x\neq1$ 时，
\[
f(x)=\frac{(x-1)(x+1)}{x-1}=x+1\text{。}
\]
所以
\[
\lim_{x\to1}f(x)=\lim_{x\to1}(x+1)=2\text{。}
\]
极限存在且为有限数，但函数在 $x=1$ 处没有定义。因此 $x=1$ 是可去间断点。若补定义 $f(1)=2$，所得函数将在 $x=1$ 处连续。\par
% UNIT S04.U10
\begin{example}
设
\[
f(x)=
\begin{cases}
1+x,&x<0,\\
2+x^2,&x\geq0\text{。}
\end{cases}
\]
判断 $f$ 在 $x=0$ 处的连续性及间断类型。
\end{example}
% UNIT S04.U11
\par\noindent\textbf{解：}由左侧表达式可得
\[
\lim_{x\to0^-}f(x)=\lim_{x\to0^-}(1+x)=1\text{。}
\]
由右侧表达式可得
\[
\lim_{x\to0^+}f(x)=\lim_{x\to0^+}(2+x^2)=2\text{。}
\]
又因为 $f(0)=2$，所以右极限等于函数值，但左极限与右极限不相等。根据双侧极限存在的充要条件，$\lim_{x\to0}f(x)$ 不存在。因此 $f$ 在 $x=0$ 处不连续，且 $x=0$ 是跳跃间断点。\par
% UNIT S04.U12
\begin{exercise}
\label{mathagent:exercise:89dc01d13bdd47d62d3f0c96d9d7bd4201cee6446e15a7269f32977f2b24d58c}
设
\[
f(x)=
\begin{cases}
\dfrac{x^2-4}{x-2},&x\neq2,\\
a,&x=2\text{。}
\end{cases}
\]
求常数 $a$ 的值，使函数 $f$ 在 $x=2$ 处连续，并说明若 $a$ 取其他值时该点的间断类型。
\end{exercise}
\section{闭区间连续函数的基本定理及应用}\label{sec:S05}
% UNIT S05.U01
本节从点附近的连续性转向整个闭区间上的结论。沿用定义 \ref{def:interval-continuity}：在 $[a,b]$ 上连续包括内点连续、左端点右连续和右端点左连续。

证明中使用实数的两条基本性质：非空且有上界的实数集合存在上确界（最小上界）；一列逐层包含的非空闭区间，若长度趋于零，则恰有一个公共点。这些性质作为实数的基础性质使用，不在此展开其完备性证明。\par
% UNIT S05.U02
\begin{definition}
\label{def:bounded-on-closed-interval}设 $a,b\in\mathbb{R}$，$a\leq b$，函数 $f:[a,b]\to\mathbb{R}$。若存在常数 $K>0$，使得对每个 $x\in[a,b]$ 都有 $|f(x)|\leq K$，则称 $f$ 在 $[a,b]$ 上有界。

若存在 $u\in[a,b]$，使得对所有 $x\in[a,b]$ 都有 $f(x)\leq f(u)$，则称 $f(u)$ 为最大值；若存在 $v\in[a,b]$，使得对所有 $x\in[a,b]$ 都有 $f(v)\leq f(x)$，则称 $f(v)$ 为最小值。有界只要求存在统一的界，取得最值还要求相应的值确实由定义域中的点取到。
\end{definition}
% UNIT S05.U03
\begin{theorem}
\label{thm:extreme-value-theorem}有界性与最值定理：设 $a,b\in\mathbb{R}$，$a\leq b$，函数 $f:[a,b]\to\mathbb{R}$ 在 $[a,b]$ 上连续，则 $f$ 在该区间上有界，并且存在 $u,v\in[a,b]$，使得
\[
f(v)\leq f(x)\leq f(u)\qquad\text{对所有 }x\in[a,b].
\]
因此，闭区间上的连续函数必能取得最大值和最小值；取得最值的点不一定唯一，也不一定是端点。
\end{theorem}
% UNIT S05.U04
\begin{proof}
若 $a=b$，定义域只有一点，结论显然成立。以下设 $a<b$。

先证明有界性。反设 $f$ 在 $[a,b]$ 上无界。把区间平分，至少有一个半区间上 $f$ 无界，否则两个半区间各自的界取较大者便是整个区间的界。反复选择这样的半区间，得到逐层包含的闭区间 $I_n$，其中 $I_0=[a,b]$，$I_n$ 的长度为 $(b-a)/2^n$，且 $f$ 在每个 $I_n$ 上均无界。

由闭区间套性质，存在公共点 $c\in[a,b]$。由 $f$ 在 $c$ 处连续（端点采用相应的单侧连续），存在 $\delta>0$，使得对 $x\in[a,b]$，只要 $|x-c|<\delta$，就有
\[
|f(x)-f(c)|<1,\qquad |f(x)|<|f(c)|+1.
\]
取充分大的正整数 $n$，使 $(b-a)/2^n<\delta$。因为 $c\in I_n$，每个 $x\in I_n$ 都满足 $|x-c|<\delta$，于是 $f$ 在 $I_n$ 上有界，矛盾。

再证明最大值存在。值域 $\{f(x):x\in[a,b]\}$ 非空且有界，令其上确界为 $M$。若没有任何点取到 $M$，则对每个 $x\in[a,b]$ 都有 $M-f(x)>0$。由连续函数的运算性质，函数
\[
g(x)=\frac{1}{M-f(x)}\qquad(x\in[a,b])
\]
在 $[a,b]$ 上连续。因此由刚证明的有界性，存在 $K>0$，使 $0<g(x)\leq K$ 对所有 $x\in[a,b]$ 成立。这意味着
\[
f(x)\leq M-\frac{1}{K}\qquad(x\in[a,b]),
\]
从而比 $M$ 更小的数 $M-1/K$ 也是值域的上界，与 $M$ 是最小上界矛盾。所以存在 $u\in[a,b]$ 使 $f(u)=M$。

将最大值的结论用于连续函数 $-f$，得到某个 $v\in[a,b]$ 使 $-f(x)\leq-f(v)$ 对所有 $x\in[a,b]$ 成立，即 $f(v)$ 是 $f$ 的最小值。
\end{proof}
% UNIT S05.U05
\begin{example}
分别考察下列函数，说明闭区间连续性定理的条件为什么不能随意省略。
\[
\begin{aligned}
&f(x)=x,&&x\in(0,1);\\
&g(x)=\frac{1}{x},&&x\in(0,1];\\
&h(x)=\begin{cases}x,&0\leq x<1,\\0,&x=1,\end{cases}&&x\in[0,1].
\end{aligned}
\]
判断它们是否有界以及是否取得最大值、最小值。
\end{example}
% UNIT S05.U06
\par\noindent\textbf{解：}$f$ 有界，但没有最大值和最小值：对任意 $x\in(0,1)$，都有 $x/2<x<(x+1)/2$，且这三个点均在定义域内。它连续，但定义域不是闭区间。

$g$ 无界，因为对任意 $K>0$，取 $x=1/(K+1)\in(0,1]$，便有 $g(x)=K+1>K$。它在 $x=1$ 处取得最小值 $1$，没有最大值；缺失端点 $0$ 的条件不能忽略。

$h$ 有界，在 $x=0$ 和 $x=1$ 处取得最小值 $0$，但没有最大值：对 $0\leq x<1$，点 $(x+1)/2$ 的函数值更大；$x=1$ 处的值也不是最大值。它虽然定义在闭区间上，却在右端点不连续，因为
\[
\lim_{x\to1^-}h(x)=1\neq0=h(1).
\]
这些例子说明，定理中的闭区间与连续性是保证结论的条件；缺少条件时，结论可能失效，而不是必然全部失效。\par
% UNIT S05.U07
\begin{theorem}
\label{thm:zero-existence-theorem}零点存在性定理：设 $a,b\in\mathbb{R}$，$a<b$，函数 $f:[a,b]\to\mathbb{R}$ 在 $[a,b]$ 上连续。若 $f(a)f(b)<0$，则至少存在一点 $c\in(a,b)$，使得 $f(c)=0$。

若某个端点的函数值已经等于零，则该端点本身就是 $[a,b]$ 中的零点，但这并不保证开区间 $(a,b)$ 内还有零点。
\end{theorem}
% UNIT S05.U08
\begin{proof}
先设 $f(a)<0<f(b)$。令 $[a_0,b_0]=[a,b]$。在已选区间 $[a_n,b_n]$ 中取中点 $m_n=(a_n+b_n)/2$。若 $f(m_n)=0$，则找到所求零点；否则，若 $f(m_n)<0$，取下一段为 $[m_n,b_n]$；若 $f(m_n)>0$，取下一段为 $[a_n,m_n]$。

若过程一直没有遇到零点，就得到逐层包含的闭区间，且对每个非负整数 $n$，
\[
f(a_n)<0<f(b_n),\qquad b_n-a_n=\frac{b-a}{2^n}.
\]
由闭区间套性质，存在公共点 $c\in[a,b]$，并且
\[
|a_n-c|\leq\frac{b-a}{2^n},\qquad |b_n-c|\leq\frac{b-a}{2^n}.
\]
若 $f(c)>0$，由连续性，存在 $\delta>0$，使得当 $x\in[a,b]$ 且 $|x-c|<\delta$ 时，有 $|f(x)-f(c)|<f(c)/2$，因而 $f(x)>0$。取充分大的 $n$，便得 $f(a_n)>0$，矛盾。若 $f(c)<0$，同理利用误差 $|f(c)|/2$ 可得充分大的 $n$ 满足 $f(b_n)<0$，也矛盾。因此 $f(c)=0$。由于两个端点的值均不为零，$c\in(a,b)$。

若 $f(a)>0>f(b)$，对连续函数 $-f$ 应用上述论证，即得同样结论。
\end{proof}
% UNIT S05.U09
\begin{theorem}
\label{thm:intermediate-value-theorem}介值定理：设 $a,b\in\mathbb{R}$，$a<b$，函数 $f:[a,b]\to\mathbb{R}$ 在 $[a,b]$ 上连续。对任意满足
\[
\min\{f(a),f(b)\}\leq y\leq\max\{f(a),f(b)\}
\]
的实数 $y$，都存在 $c\in[a,b]$，使得 $f(c)=y$。若 $y$ 严格介于两个端点函数值之间，则可以取 $c\in(a,b)$。

结合定理 \ref{thm:extreme-value-theorem}，若 $m,M$ 分别为 $f$ 在 $[a,b]$ 上的最小值与最大值，则 $f$ 的值域恰为 $[m,M]$。
\end{theorem}
% UNIT S05.U10
\begin{proof}
若 $y=f(a)$ 或 $y=f(b)$，取相应端点即可。若 $y$ 严格介于 $f(a)$ 与 $f(b)$ 之间，令
\[
F(x)=f(x)-y\qquad(x\in[a,b]).
\]
由连续函数的运算性质，$F$ 在 $[a,b]$ 上连续，且 $F(a)F(b)<0$。由定理 \ref{thm:zero-existence-theorem}，存在 $c\in(a,b)$ 使 $F(c)=0$，即 $f(c)=y$。

最后，最值定理保证存在 $u,v\in[a,b]$，使 $f(u)=M$，$f(v)=m$。若 $m=M$，函数恒等于该常数，值域就是单点区间 $[m,M]$。若 $m<M$，则 $u\neq v$；在闭区间 $[\min\{u,v\},\max\{u,v\}]$ 上应用刚证明的介值结论，知每个 $y\in[m,M]$ 都能取到。另一方面，由最值定义，每个函数值都在 $[m,M]$ 内，故值域恰为 $[m,M]$。
\end{proof}
% UNIT S05.U11
\begin{example}
设 $f(x)=x^2-4x+5$，$x\in[0,3]$。不用导数，求最大值、最小值及值域，并判断方程 $f(x)=3$ 在此区间内是否有解。
\end{example}
% UNIT S05.U12
\par\noindent\textbf{解：}多项式连续，故最值定理适用。对所有实数 $x$，配方恒等式为
\[
x^2-4x+5=(x-2)^2+1.
\]
对 $x\in[0,3]$，有 $-2\leq x-2\leq1$，故 $0\leq(x-2)^2\leq4$。下界在 $x=2$ 处取到，上界在 $x=0$ 处取到，所以最小值为 $1$，最大值为 $5$。由介值定理，值域为 $[1,5]$，因而 $f(x)=3$ 至少有一解。

具体解方程得 $x=2\pm\sqrt{2}$。由于 $1<\sqrt{2}<2$，只有 $2-\sqrt{2}$ 属于 $[0,3]$，所以此区间内恰有一个解。这里的唯一性来自具体代数求解，而不是介值定理。\par
% UNIT S05.U13
\begin{example}
考虑某商品的简化教学模型：$p$ 为以元/千克计的价格数值，$D(p)$ 与 $S(p)$ 分别为同一个固定的一周内、以千克计的需求量和供给量数值。下列公式是这些统一单位下数值之间的假设关系，并非对真实市场的事实描述。在价格范围 $p\in[0,8]$（即每千克 $0$ 至 $8$ 元）内，某商品的需求量与供给量分别为
\[
D(p)=24-2p,\qquad S(p)=4+p.
\]
先用连续性说明该价格范围内存在供需相等的均衡价格，再求出均衡价格与交易量，并说明唯一性的依据。
\end{example}
% UNIT S05.U14
\par\noindent\textbf{解：}沿用题面的简化教学模型及单位：价格以元/千克计，供需数量均为同一个固定的一周内以千克计的数值。在 $[0,8]$ 上，两种数量均非负，且 $D,S$ 均为连续函数。定义超额需求
\[
E(p)=D(p)-S(p)=20-3p\qquad(p\in[0,8]).
\]
这里 $E(p)$ 也是该周内以千克计的数量差数值，价格区间 $[0,8]$ 对应每千克 $0$ 至 $8$ 元。它在闭区间上连续，且 $E(0)=20>0$，$E(8)=-4<0$，分别表示该周的需求量比供给量多 $20$ 千克和少 $4$ 千克。由零点存在性定理，存在 $p_*\in(0,8)$ 使 $E(p_*)=0$，即供需相等。

解线性方程 $20-3p_*=0$，得
\[
p_*=\frac{20}{3},\qquad D(p_*)=S(p_*)=\frac{32}{3}.
\]
因此，在此教学模型中，均衡价格为 $\frac{20}{3}$ 元/千克，均衡交易量为该周内 $\frac{32}{3}$ 千克。
对任意 $0\leq p_1<p_2\leq8$，都有
\[
E(p_2)-E(p_1)=-3(p_2-p_1)<0.
\]
因此超额需求严格递减，不可能有两个不同零点。存在性来自连续性与端点异号，唯一性则需要这里的严格递减性质；一般连续模型未必只有一个均衡点。\par
% UNIT S05.U15
应用时应依次检查：定义域是否为有限闭区间；函数是否在整个区间上连续，包括端点；所求目标值是否介于两个已知函数值之间。对方程 $f(x)=g(x)$，通常先令 $F(x)=f(x)-g(x)$，再检查连续性与异号条件。

零点定理给出的是充分条件：端点同号不等于没有零点；端点异号也不等于只有一个零点。最值定理保证最值能取到，但不直接给出最值的位置或数值。以下练习的答案集中放在本节最后，作为讲义末尾的练习答案。\par
% UNIT S05.U16
\begin{exercise}
\label{mathagent:exercise:8df9c883175539781c5297e80328cb621961e5226b4ee75741c01793ba96210d}
不用导数，求函数 $f(x)=x^2-2x+3$ 在 $[0,3]$ 上的最小值、最大值及值域。再判断方程 $f(x)=4$ 在该区间内有几个解。
\end{exercise}
% UNIT S05.U17
\begin{exercise}
\label{mathagent:exercise:15f24ef1d79299958d598c18e2c5f872d77ddb70db01666a2348f05d38173ea8}
设 $F(x)=x^3+x-1$，$x\in[0,1]$。证明方程 $F(x)=0$ 在 $(0,1)$ 内至少有一个解，并用因式分解说明这个解唯一。
\end{exercise}
% UNIT S05.U18
\begin{exercise}
\label{mathagent:exercise:14eff0c9afb4ff01b5d129f61082fbaf47f00da9185eedc3ed8874ad925b6e49}
考虑某商品的简化教学模型：$p$ 为以元/千克计的价格数值，需求函数 $D$ 和供给函数 $S$ 的值分别为同一个固定的一周内、以千克计的需求量和供给量数值。以下条件是统一单位下数值之间的教学假设，并非对真实市场的事实描述。某商品在 $p\in[1,5]$（即每千克 $1$ 至 $5$ 元）上的需求函数 $D$ 和供给函数 $S$ 均连续且非负，满足
\[
D(1)=12,\quad S(1)=3,\qquad D(5)=4,\quad S(5)=8.
\]
证明存在区间内部的均衡价格。若进一步知道 $D$ 严格递减、$S$ 严格递增，证明均衡价格唯一。仅凭最初给出的条件能否断言唯一？
\end{exercise}
% UNIT S05.U19
\textbf{练习答案（讲义末尾）}

以下按上述练习顺序给出答案与必要说明。\par
\section*{练习答案}
% UNIT S01.U13
\par\noindent\textbf{\hyperref[mathagent:exercise:bf60311c135cf0c14f3cb42069a931ed175758be7d89dff861bed83a8d332803]{练习 \ref*{mathagent:exercise:bf60311c135cf0c14f3cb42069a931ed175758be7d89dff861bed83a8d332803} 的解答：}}当 $x\to1^-$ 时，$f(x)=2x+1$，因此
\[
\lim_{x\to1^-}f(x)=2\cdot1+1=3\text{。}
\]
当 $x\to1^+$ 时，$f(x)=5$，因此
\[
\lim_{x\to1^+}f(x)=5\text{。}
\]
由于左极限 $3$ 与右极限 $5$ 不相等，根据定理 \ref{thm:two-sided-limit-criterion}，双侧极限 $\lim_{x\to1}f(x)$ 不存在。\par
% UNIT S02.U13
\par\noindent\textbf{\hyperref[mathagent:exercise:b4c9c9a1fb82bcef052fdbf6e1539831ed88cee016dbcbf83a50e5e8d2311be0]{练习 \ref*{mathagent:exercise:b4c9c9a1fb82bcef052fdbf6e1539831ed88cee016dbcbf83a50e5e8d2311be0} 的解答：}}由极限的和与差法则，
\[
\lim_{x\to x_0}(f(x)+g(x))=2+(-1)=1,
\]
\[
\lim_{x\to x_0}(f(x)-g(x))=2-(-1)=3.
\]
分母的极限为 $3\neq0$，所以分母在 $x_0$ 的某个去心邻域内不为零，可以使用商的极限法则。因此
\[
\lim_{x\to x_0}\frac{f(x)+g(x)}{f(x)-g(x)}=\frac{1}{3}.
\]

附加小问解答：在去心邻域 $0<|x|<1$ 内，$1/x$ 有定义。由余弦函数的有界性，
\[
-1\leq\cos(1/x)\leq1,
\qquad
1\leq2+\cos(1/x)\leq3.
\]
因此分母始终为正，原式在该去心邻域内有定义。对上述正数取倒数，再乘以非负数 $x^2$，得到
\[
\frac{x^2}{3}\leq\frac{x^2}{2+\cos(1/x)}\leq x^2.
\]
上下界也在同一去心邻域内有定义，三者均考察双侧趋近过程 $x\to0$。由 $\lim_{x\to0}x=0$ 及极限的积与常数倍法则，
\[
\lim_{x\to0}\frac{x^2}{3}=0,
\qquad
\lim_{x\to0}x^2=0.
\]
两端趋于同一有限数，故由夹逼定理，
\[
\lim_{x\to0}\frac{x^2}{2+\cos(1/x)}=0.
\]
这里仅使用余弦函数的有界性与代数极限，不需要求 $\cos(1/x)$ 的极限，也不需要在 $x=0$ 处定义原式。\par
% UNIT S03.U10
\par\noindent\textbf{\hyperref[mathagent:exercise:410f498083ff9a28d1a2acf9839e74554b434d3445561642fdfe5d3854f06971]{练习 \ref*{mathagent:exercise:410f498083ff9a28d1a2acf9839e74554b434d3445561642fdfe5d3854f06971} 的解答：}}当 $x\to0$ 且 $x\neq0$ 时，$\dfrac{e^x-1}{x}$ 有定义。由基本极限
\[
\lim_{x\to0}\frac{e^x-1}{x}=1,
\]
以及
\[
\lim_{x\to0}(1-2x)=1,
\]
再由极限的乘法法则，得到
\[
\lim_{x\to0}\frac{e^x-1}{x}(1-2x)=1\cdot1=1.
\]
乘法法则要求两个因子在同一去心邻域内有定义，并且各自的有限极限存在；这些条件在本题中均满足。\par
% UNIT S04.U13
\par\noindent\textbf{\hyperref[mathagent:exercise:89dc01d13bdd47d62d3f0c96d9d7bd4201cee6446e15a7269f32977f2b24d58c]{练习 \ref*{mathagent:exercise:89dc01d13bdd47d62d3f0c96d9d7bd4201cee6446e15a7269f32977f2b24d58c} 的解答：}}当 $x\neq2$ 时，
\[
\frac{x^2-4}{x-2}=\frac{(x-2)(x+2)}{x-2}=x+2\text{。}
\]
因此
\[
\lim_{x\to2}f(x)=\lim_{x\to2}(x+2)=4\text{。}
\]
函数在 $x=2$ 处连续的条件是
\[
\lim_{x\to2}f(x)=f(2)=a\text{，}
\]
所以必须取
\[
a=4\text{。}
\]
当 $a=4$ 时，函数在 $x=2$ 处连续；当 $a\neq4$ 时，极限仍存在且等于 $4$，但函数值不等于极限值，因此 $x=2$ 是可去间断点。\par
% UNIT S05.U20
\par\noindent\textbf{\hyperref[mathagent:exercise:8df9c883175539781c5297e80328cb621961e5226b4ee75741c01793ba96210d]{练习 \ref*{mathagent:exercise:8df9c883175539781c5297e80328cb621961e5226b4ee75741c01793ba96210d} 的解答：}}配方得 $f(x)=(x-1)^2+2$。当 $x\in[0,3]$ 时，$-1\leq x-1\leq2$，故 $0\leq(x-1)^2\leq4$。最小值为 $f(1)=2$，最大值为 $f(3)=6$。由连续性与介值定理，值域为 $[2,6]$。

方程 $f(x)=4$ 等价于 $(x-1)^2=2$，实数解为 $1\pm\sqrt{2}$。其中 $1-\sqrt{2}<0$，而 $1+\sqrt{2}\in[0,3]$，故区间内恰有一个解。\par
% UNIT S05.U21
\par\noindent\textbf{\hyperref[mathagent:exercise:15f24ef1d79299958d598c18e2c5f872d77ddb70db01666a2348f05d38173ea8]{练习 \ref*{mathagent:exercise:15f24ef1d79299958d598c18e2c5f872d77ddb70db01666a2348f05d38173ea8} 的解答：}}$F$ 为多项式，在 $[0,1]$ 上连续，且 $F(0)=-1<0$，$F(1)=1>0$，所以由零点存在性定理，存在 $c\in(0,1)$ 使 $F(c)=0$。

若 $0\leq x_1<x_2\leq1$，则
\[
F(x_2)-F(x_1)=(x_2-x_1)(x_2^2+x_1x_2+x_1^2+1)>0.
\]
因此 $F$ 严格递增，若有两个不同零点，就会与此严格不等式矛盾。故零点唯一。\par
% UNIT S05.U22
\par\noindent\textbf{\hyperref[mathagent:exercise:14eff0c9afb4ff01b5d129f61082fbaf47f00da9185eedc3ed8874ad925b6e49]{练习 \ref*{mathagent:exercise:14eff0c9afb4ff01b5d129f61082fbaf47f00da9185eedc3ed8874ad925b6e49} 的解答：}}沿用题面的简化教学模型，价格以元/千克计，需求量与供给量均指同一个固定的一周内以千克计的数量。令 $E(p)=D(p)-S(p)$，$p\in[1,5]$，其中 $E(p)$ 为该周超额需求的千克数，价格区间对应每千克 $1$ 至 $5$ 元。由连续函数的差仍连续，$E$ 在 $[1,5]$ 上连续，且
\[
E(1)=9>0,\qquad E(5)=-4<0.
\]
这分别表示该周的需求量比供给量多 $9$ 千克和少 $4$ 千克。所以存在 $p_*\in(1,5)$ 使 $E(p_*)=0$，即 $D(p_*)=S(p_*)$。均衡价格为 $p_*$ 元/千克，位于每千克 $1$ 元与 $5$ 元之间；相等的供需数量是该周以千克计的均衡交易量。

若进一步有需求严格递减、供给严格递增，则对任意 $1\leq p_1<p_2\leq5$，
\[
E(p_1)-E(p_2)=[D(p_1)-D(p_2)]+[S(p_2)-S(p_1)]>0.
\]
因此 $E$ 严格递减，零点唯一。

仅凭最初的连续性和端点数据不能断言唯一。例如，令 $D,S$ 分别在相邻整数价格之间作直线连接，其在 $p=1,2,3,4,5$ 处的值依次为
\[
\begin{array}{c|ccccc}
p&1&2&3&4&5\\
D(p)&12&4&8&4&4\\
S(p)&3&5&5&5&8
\end{array}
\]
表中第一行是以元/千克计的价格数值，后两行是该固定一周内以千克计的供需数量数值；此表仅为教学模型中的反例数据。所得函数连续且非负，并满足给定端点条件。

超额需求在这五个价格处的值依次为
\[
9,-1,3,-1,-4.
\]
这些数值同样以该周的千克数计。因此在 $(1,2)$、$(2,3)$、$(3,4)$ 各有一个零点，说明原条件允许多个均衡价格；这三个开区间分别对应每千克 $1$ 至 $2$ 元、$2$ 至 $3$ 元、$3$ 至 $4$ 元之间的价格范围。\par
\end{document}
